\documentclass[letterpaper,12pt,twoside,]{pinp}

%% Some pieces required from the pandoc template
\providecommand{\tightlist}{%
  \setlength{\itemsep}{0pt}\setlength{\parskip}{0pt}}

% Use the lineno option to display guide line numbers if required.
% Note that the use of elements such as single-column equations
% may affect the guide line number alignment.

\usepackage[T1]{fontenc}
\usepackage[utf8]{inputenc}
\usepackage{longtable}

% pinp change: the geometry package layout settings need to be set here, not in pinp.cls
\geometry{layoutsize={0.95588\paperwidth,0.98864\paperheight},%
  layouthoffset=0.02206\paperwidth, layoutvoffset=0.00568\paperheight}

\definecolor{pinpblue}{HTML}{185FAF}  % imagecolorpicker on blue for new R logo
\definecolor{pnasbluetext}{RGB}{101,0,0} %


\usepackage{booktabs}
\usepackage{longtable}
\usepackage{array}
\usepackage{multirow}
\usepackage{wrapfig}
\usepackage{float}
\usepackage{colortbl}
\usepackage{pdflscape}
\usepackage{tabu}
\usepackage{threeparttable}
\usepackage{threeparttablex}
\usepackage[normalem]{ulem}
\usepackage{makecell}

\title{Assignment 2 - Central Limit Theorem, Confidence Intervals and
Bootstrap. Due October 11, 11:59pm 2019}

\author[a]{EPIB607 - Inferential Statistics}

  \affil[a]{Fall 2020, McGill University}

\setcounter{secnumdepth}{5}

% Please give the surname of the lead author for the running footer
\leadauthor{Bhatnagar}

% Keywords are not mandatory, but authors are strongly encouraged to provide them. If provided, please include two to five keywords, separated by the pipe symbol, e.g:
 \keywords{  Sampling distribution |  Standard error |  Normal distribution |  Quantiles |  Percentiles |  Z-scores  }  

\begin{abstract}
In this assignment you will practice calculating confidence intervals.
Provide confidence intervals with appropriate units. Answers should be
given in full sentences (DO NOT just provide the number). All graphs and
calculations are to be completed in an R Markdown document using the
provided template. You are free to choose any function from any package
to complete the assignment. Concise answers will be rewarded. Be brief
and to the point. Please submit the compiled pdf report to Crowdmark by
October 11, 2020, 11:59pm. You need to save your answers to each
question in separate pdf files. You also need to upload your code. See
\url{https://crowdmark.com/help/} for details.
\end{abstract}

\dates{This version was compiled on \today} 

% initially we use doi so keep for backwards compatibility
% new name is doi_footer

\pinpfootercontents{Assignment 2 due October 11, 2020 by 11:59pm}

\begin{document}

% Optional adjustment to line up main text (after abstract) of first page with line numbers, when using both lineno and twocolumn options.
% You should only change this length when you've finalised the article contents.
\verticaladjustment{-2pt}

\maketitle
\thispagestyle{firststyle}
\ifthenelse{\boolean{shortarticle}}{\ifthenelse{\boolean{singlecolumn}}{\abscontentformatted}{\abscontent}}{}

% If your first paragraph (i.e. with the \dropcap) contains a list environment (quote, quotation, theorem, definition, enumerate, itemize...), the line after the list may have some extra indentation. If this is the case, add \parshape=0 to the end of the list environment.


\hypertarget{points-2-each-rainfall}{%
\section{(8 points, 2 each) Rainfall}\label{points-2-each-rainfall}}

Statistics from Cornell's Northeast Regional Climate Center indicate
that Ithaca, NY, gets an average of 35.4 feet of rain each year, with a
standard deviation of 4.2 feet. Assume that a Normal model applies.

\begin{enumerate}
\def\labelenumi{\alph{enumi}.}
\tightlist
\item
  During what percentage of years does Ithaca get more than 40 feet of
  rain?
\item
  Less than how much rain falls in the driest 20\% of all years?
\item
  A Cornell University student is in Ithaca for 4 years. Let \(\bar{y}\)
  represent the mean amount of rain for those 4 years. Describe the
  sampling distribution model of this sample mean, \(\bar{y}\).
\item
  What's the probability that those 4 years average less than 30 feet of
  rain?
\end{enumerate}

\hypertarget{points-museum-membership}{%
\section{(12 points) Museum membership}\label{points-museum-membership}}

A museum offers several levels of membership, as shown in the following
table:

\begin{longtable}[]{@{}lll@{}}
\toprule
Member Category & Amount of Donation (\$) & Percent of
Members\tabularnewline
\midrule
\endhead
Individual & 50 & 41\tabularnewline
Family & 100 & 37\tabularnewline
Sponsor & 250 & 14\tabularnewline
Patron & 500 & 7\tabularnewline
Benefactor & 1000 & 1\tabularnewline
\bottomrule
\end{longtable}

\begin{enumerate}
\def\labelenumi{\alph{enumi}.}
\tightlist
\item
  (2 points) Find the mean and standard deviation of the donations.
\item
  (4 points) During their annual membership drive, they hope to sign up
  50 new members each day. Would you expect the distribution of the
  donations for a day to follow a Normal model? Explain.
\item
  (2 points) Consider the mean donation of the 50 new members each day.
  Describe the sampling model for these means (shape, center, and
  spread).
\item
  (4 points) Assume a hypothetical cohort of 1000 members follows the
  distribution in the Table. Plot the frequency histogram of the
  donation amounts.
\end{enumerate}

\hypertarget{points-where-do-you-buy}{%
\section{(12 points) Where do you buy?}\label{points-where-do-you-buy}}

Consumers can purchase nonprescription medications at food stores, mass
merchandise stores such as Kmart and Wal-Mart, or pharmacies. About 45\%
of consumers make such purchases at pharmacies. What accounts for the
popularity of pharmacies, which often charge higher prices?

A study examined consumers' perceptions of overall performance of the
three types of stores, using a long questionnaire that had asked about
such things as ``neat and attractive store'', ``knowledgeable staff'',
and ``assistance in choosing among various types of nonprescription
medication''. A performance score was based on 27 such questions. The
subjects were 201 people chosen at random from the Indianapollis
telephone directory. Here are the means (\(\bar{y}\)) and standard
deviations (\(s\)) of the performances scores for the sample:

\begin{longtable}[]{@{}lll@{}}
\toprule
Store type & \(\bar{y}\) & \(s\)\tabularnewline
\midrule
\endhead
Food stores & 18.67 & 24.95\tabularnewline
Mass merchandisers & 32.38 & 33.37\tabularnewline
Pharmacies & 48.6 & 35.62\tabularnewline
\bottomrule
\end{longtable}

\begin{enumerate}
\def\labelenumi{\alph{enumi}.}
\tightlist
\item
  (2 points) What population do you think the authors of the study want
  to draw conclusions about?
\item
  (2 points) What population are you certain they can draw conclusions
  about?
\item
  (5 points) Give a 95\% confidence interval for the mean performance
  for each type of store.
\item
  (3 points) Based on these confidence intervals, are you convinced that
  consumers think that pharmacies offer higher performance than the
  other types of stores?
\end{enumerate}

\hypertarget{points-deer-mice}{%
\section{(18 points) Deer mice}\label{points-deer-mice}}

Deer mice are small rodents native to North America. Their adult body
lengths (excluding tail) are known to vary approximately Normally, with
mean \(\mu = 86\) millimeters (mm) and standard deviation \(\sigma=8\)
mm. Deer mice are found in diverse habitats and exhibit different
adaptations to their environment. A random sample of 14 deer mice in a
rich forest habitat gives an average body length of \(\bar{y} = 91.1\)
mm. Assume that the standard deviation \(\sigma\) of all deer mice in
this area is also 8 mm.

\begin{enumerate}
\def\labelenumi{\alph{enumi}.}
\tightlist
\item
  (3 points) What is the standard deviation of the mean length
  \(\bar{y}\)?
\item
  (5 points) What critical value do you need to use in order to compute
  a 95\% confidence interval for the mean \(\mu\)
\item
  (7 points) Give a i) 90\% confidence interval and ii) 95\% confidence
  interval for the mean body length of all deer mice in the forest
  habitat.
\item
  (3 points) Why does the 90\% interval have a smaller margin of error?
\end{enumerate}

\hypertarget{points-does-breast-feeding-weaken-bones}{%
\section{(20 points) Does breast-feeding weaken
bones?}\label{points-does-breast-feeding-weaken-bones}}

Breast-feeding mothers secrete calcium into their milk. Some of the
calcium may come from their bones, so mothers may lose bone mineral
content. Researchers compared 47 breast-feeding women with 22 women of
similar age who were neither pregnant nor lactating. They measured the
percent change in the mineral content of the women's spines over three
months. A negative value of \texttt{mineral\_loss} indicates a loss in
mineral content. The data can be read into \texttt{R} and a boxplot
comparing the two groups is shown in the Figure below:

\begin{Shaded}
\begin{Highlighting}[]
\NormalTok{bnurl <-}\StringTok{ "https://github.com/sahirbhatnagar/EPIB607/raw/master/data/boneloss.csv"}
\NormalTok{boneloss <-}\StringTok{ }\KeywordTok{read.csv}\NormalTok{(bnurl)}
\KeywordTok{head}\NormalTok{(boneloss)}
\end{Highlighting}
\end{Shaded}

\begin{ShadedResult}
\begin{verbatim}
#     type mineral_loss
#  1 Other          2.4
#  2 Other          0.0
#  3 Other          0.9
#  4 Other         -0.2
#  5 Other          1.0
#  6 Other          1.7
\end{verbatim}
\end{ShadedResult}

\begin{Shaded}
\begin{Highlighting}[]
\KeywordTok{library}\NormalTok{(mosaic)}
\NormalTok{ggformula}\OperatorTok{::}\KeywordTok{gf_boxplot}\NormalTok{(mineral_loss }\OperatorTok{~}\StringTok{ }\NormalTok{type, }\DataTypeTok{data =}\NormalTok{ boneloss)}
\end{Highlighting}
\end{Shaded}

\begin{center}\includegraphics{a2-clt-ci_files/figure-latex/unnamed-chunk-3-1} \end{center}

\begin{enumerate}
\def\labelenumi{\alph{enumi}.}
\tightlist
\item
  (4 points) Give a 95\% confidence interval for the mean mineral loss
  for each group.
\item
  (6 points) Based on these confidence intervals, are you convinced that
  the data show distinctly greater bone mineral loss among the
  breast-feeding women?
\item
  (10 points) Let \(\bar{y}_1\) be the mean mineral loss for
  breast-feeding women and \(\bar{y}_2\) be the mean mineral loss for
  the other group. Construct a 95\% bootstrap confidence interval for
  the difference \(\bar{y}_2 - \bar{y}_1\) based on \(B=10000\)
  bootstrap samples. Interpret this confidence interval and compare it
  to the ones obtained in b. \emph{Hint}: split the data in two
  data.frames by \texttt{type}. For each data.frame, create \(B\)
  resamples and calculate the means for each. Plot the histogram of the
  difference in means, which gives you the bootstrap distribution for
  the difference. From this you can use the \texttt{quantile} function
  to calculate the 95\% CI.
\end{enumerate}

\hypertarget{points-how-deep-is-the-ocean}{%
\section{(15 points) How deep is the
ocean?}\label{points-how-deep-is-the-ocean}}

This question is based on the
\href{https://github.com/sahirbhatnagar/EPIB607/blob/master/labs/003-ocean-depths/water_exercise_epib607.pdf}{in-class
Exercise} on sampling distributions. For your sample of \(n=20\) of
depths of the ocean, calculate the

\begin{enumerate}
\def\labelenumi{\alph{enumi}.}
\tightlist
\item
  (3 points) 95\% Confidence interval using the \(\pm\) formula
\item
  (3 points) 95\% Confidence interval using the \texttt{qnorm} function
\item
  (5 points) 95\% Confidence interval using \(B=10000\) bootstrap
  samples
\item
  (4 points) Plot all three confidence intervals on the same plot and
  comment on the difference/similarities between the 3 intervals. You
  may use the \texttt{compare\_CI} function provided below to produce
  the plot. This takes as input, the sample mean (\texttt{ybar}), and
  the CIs calculated from a,b,c in the form of a numeric vector of
  length 2 into the arguments \texttt{PM}, \texttt{QNORM} and
  \texttt{BOOT}, respectively.
\end{enumerate}

\begin{Shaded}
\begin{Highlighting}[]
\NormalTok{compare_CI <-}\StringTok{ }\ControlFlowTok{function}\NormalTok{(ybar, PM, QNORM, BOOT, }
                       \DataTypeTok{col =} \KeywordTok{c}\NormalTok{(}\StringTok{"#E41A1C"}\NormalTok{,}\StringTok{"#377EB8"}\NormalTok{,}\StringTok{"#4DAF4A"}\NormalTok{)) \{}

\NormalTok{  dt <-}\StringTok{ }\KeywordTok{data.frame}\NormalTok{(}\DataTypeTok{type =} \KeywordTok{c}\NormalTok{(}\StringTok{"plus_minus"}\NormalTok{, }\StringTok{"qnorm"}\NormalTok{, }\StringTok{"bootstrap"}\NormalTok{),}
                   \DataTypeTok{ybar =} \KeywordTok{rep}\NormalTok{(ybar, }\DecValTok{3}\NormalTok{),}
                   \DataTypeTok{low =} \KeywordTok{c}\NormalTok{(PM[}\DecValTok{1}\NormalTok{], QNORM[}\DecValTok{1}\NormalTok{], BOOT[}\DecValTok{1}\NormalTok{]),}
                   \DataTypeTok{up =} \KeywordTok{c}\NormalTok{(PM[}\DecValTok{2}\NormalTok{], QNORM[}\DecValTok{2}\NormalTok{], BOOT[}\DecValTok{2}\NormalTok{])}
\NormalTok{  )}
  
  \KeywordTok{plot}\NormalTok{(dt}\OperatorTok{$}\NormalTok{ybar, }\DecValTok{1}\OperatorTok{:}\KeywordTok{nrow}\NormalTok{(dt), }\DataTypeTok{pch =} \DecValTok{20}\NormalTok{, }\DataTypeTok{ylim =} \KeywordTok{c}\NormalTok{(}\DecValTok{0}\NormalTok{, }\DecValTok{5}\NormalTok{), }
       \DataTypeTok{xlim =} \KeywordTok{range}\NormalTok{(}\KeywordTok{pretty}\NormalTok{(}\KeywordTok{c}\NormalTok{(dt}\OperatorTok{$}\NormalTok{low, dt}\OperatorTok{$}\NormalTok{up))),}
       \DataTypeTok{xlab =} \StringTok{"Depth of ocean (m)"}\NormalTok{, }\DataTypeTok{ylab =} \StringTok{"Confidence Interval Type"}\NormalTok{,}
       \DataTypeTok{las =} \DecValTok{1}\NormalTok{, }\DataTypeTok{cex.axis =} \FloatTok{0.8}\NormalTok{, }\DataTypeTok{cex =} \DecValTok{3}\NormalTok{)}
  
  \KeywordTok{abline}\NormalTok{(}\DataTypeTok{v =} \DecValTok{37}\NormalTok{, }\DataTypeTok{lty =} \DecValTok{2}\NormalTok{, }\DataTypeTok{col =} \StringTok{"black"}\NormalTok{, }\DataTypeTok{lwd =} \DecValTok{2}\NormalTok{)}
  \KeywordTok{segments}\NormalTok{(}\DataTypeTok{x0 =}\NormalTok{ dt}\OperatorTok{$}\NormalTok{low, }\DataTypeTok{x1 =}\NormalTok{ dt}\OperatorTok{$}\NormalTok{up,}
           \DataTypeTok{y0 =} \DecValTok{1}\OperatorTok{:}\KeywordTok{nrow}\NormalTok{(dt), }\DataTypeTok{lend =} \DecValTok{1}\NormalTok{,}
           \DataTypeTok{col =}\NormalTok{ col, }\DataTypeTok{lwd =} \DecValTok{4}\NormalTok{)}
  
  \KeywordTok{legend}\NormalTok{(}\StringTok{"topleft"}\NormalTok{,}
         \DataTypeTok{legend =} \KeywordTok{c}\NormalTok{(}\KeywordTok{eval}\NormalTok{(}\KeywordTok{substitute}\NormalTok{( }\KeywordTok{expression}\NormalTok{(}\KeywordTok{paste}\NormalTok{(mu,}\StringTok{" = "}\NormalTok{,}\DecValTok{37}\NormalTok{)))),}
                    \KeywordTok{sprintf}\NormalTok{(}\StringTok{"plus/minus CI: [%.f, %.f]"}\NormalTok{,PM[}\DecValTok{1}\NormalTok{], PM[}\DecValTok{2}\NormalTok{]),}
                    \KeywordTok{sprintf}\NormalTok{(}\StringTok{"qnorm CI: [%.f, %.f]"}\NormalTok{,QNORM[}\DecValTok{1}\NormalTok{], QNORM[}\DecValTok{2}\NormalTok{]),}
                    \KeywordTok{sprintf}\NormalTok{(}\StringTok{"bootstrap CI: [%.f, %.f]"}\NormalTok{,BOOT[}\DecValTok{1}\NormalTok{], BOOT[}\DecValTok{2}\NormalTok{])),}
         \DataTypeTok{lty =} \KeywordTok{c}\NormalTok{(}\DecValTok{1}\NormalTok{, }\DecValTok{1}\NormalTok{,}\DecValTok{1}\NormalTok{,}\DecValTok{1}\NormalTok{),}
         \DataTypeTok{col =} \KeywordTok{c}\NormalTok{(}\StringTok{"black"}\NormalTok{,col), }\DataTypeTok{lwd =} \DecValTok{4}\NormalTok{)}
\NormalTok{\}}

\CommentTok{# example of how to use the function:}
\KeywordTok{compare_CI}\NormalTok{(}\DataTypeTok{ybar =} \DecValTok{36}\NormalTok{, }\DataTypeTok{PM =} \KeywordTok{c}\NormalTok{(}\DecValTok{28}\NormalTok{, }\DecValTok{40}\NormalTok{), }\DataTypeTok{QNORM =} \KeywordTok{c}\NormalTok{(}\DecValTok{25}\NormalTok{,}\DecValTok{40}\NormalTok{), }\DataTypeTok{BOOT =} \KeywordTok{c}\NormalTok{(}\DecValTok{31}\NormalTok{, }\DecValTok{38}\NormalTok{))}
\end{Highlighting}
\end{Shaded}

%\showmatmethods


\bibliography{pinp}
\bibliographystyle{jss}



\end{document}

